%!TEX root = ./main.tex

\section[OFDM]{The Waveform We Already Have}

% SECTION TIMING: 32:00--43:00 (6 frames). Question 2 of the roadmap frame.
% The four "Step" frames are the OFDM primer the author asked for on 2026-09-13:
% students know little OFDM, so walk one worked example (bits -> 16-QAM points ->
% perpendicular axes -> the FFT matrix as those axes -> IFFT builds the symbol and
% the grid). Keep it at the level of a 2x2 and a 4x4 matrix; no formulas beyond that.
% The OTFS / ODDM / delay-Doppler frames that used to follow were removed on 2026-09-13
% (CS audience, no signal-processing background); see parked/03_otfs_oddm_frames.tex.

% Added 2026-09-13 after the author pointed out that Step 1 began with no lead-in:
% OFDM had never been named on a slide. This frame closes question 1, poses question 2,
% answers it, and announces the four steps that follow.
\begin{frame}{Question 2: which signal do we actually use?}
  \vspace{0.3em}
  {\small Question 1 is answered: the two jobs want opposite things, so there is a
   curve of compromises. Whatever signal we pick has to carry \textbf{both} jobs and
   sit somewhere on that curve.\par}

  \vspace{0.45em}
  {\large\textbf{The answer is OFDM} \expand{Orthogonal Frequency Division
   Multiplexing} --- the signal Wi-Fi, 4G and 5G already transmit. No new radio.\par}

  \vspace{0.45em}
  {\small Before we can read OFDM as a radar we need to know what it actually is.
   Most of you have only heard the name, so we build one from scratch, in four
   steps:\par}
  \vspace{0.2em}
  {\small
  \begin{enumerate}\itemsep0pt
    \item Bits become points.
    \item Describing a point needs perpendicular axes.
    \item The FFT is a set of perpendicular axes.
    \item An OFDM symbol is one IFFT --- and symbols make the grid.
  \end{enumerate}}

  \vspace{0.15em}
  \takeaway{Four steps, one example: by the end you have built the grid this lecture is about.}
\end{frame}

\begin{frame}{Step 1: bits become points}
  \vspace{0.2em}
  \begin{columns}[T,onlytextwidth]
    \column{0.50\textwidth}
      {\small Say we want to send these bits:\par}
      \vspace{0.35em}
      {\centering\ttfamily\small 1011\ \ 0010\ \ 1110\ \ 0001\ \ \ldots\par}
      \vspace{0.7em}
      {\small Cut them into groups of four. Each group picks \textbf{one of 16 points}
       on the map at right. That step is called \textbf{16-QAM}.\par}
      \vspace{0.6em}
      {\small A point is just two numbers, $(x, y)$. So four bits became one point,
       and a long bit string became a \textbf{list of points}.\par}
    \column{0.46\textwidth}
      \centering
      \begin{tikzpicture}[scale=0.62,transform shape]
        \draw[-{Latex[length=2mm]},softgray] (-2.7,0) -- (2.7,0) node[anchor=west,font=\scriptsize]{$x$};
        \draw[-{Latex[length=2mm]},softgray] (0,-2.7) -- (0,2.7) node[anchor=south,font=\scriptsize]{$y$};
        \foreach \x in {-1.5,-0.5,0.5,1.5}
          \foreach \y in {-1.5,-0.5,0.5,1.5}
            \fill[ranblue!60] (\x,\y) circle (2.6pt);
        \fill[coreorange] (-1.5,1.5) circle (4.4pt);
        \node[font=\scriptsize\bfseries,text=coreorange,anchor=south west] at (-1.38,1.62) {1011};
        \node[font=\scriptsize,text=softgray,anchor=north] at (0,-2.85) {16 points, 4 bits each};
      \end{tikzpicture}
  \end{columns}

  \vspace{0.3em}
  \takeaway{Bits $\rightarrow$ points. Nothing has been transmitted yet;\\we have only chosen what to say.}
\end{frame}

\begin{frame}{Step 2: describing a point needs perpendicular axes}
  \vspace{0.2em}
  \begin{columns}[T,onlytextwidth]
    \column{0.48\textwidth}
      \centering
      \begin{tikzpicture}[scale=0.72,transform shape]
        \draw[-{Latex[length=2mm]},ranblue,thick] (0,0) -- (3.2,0) node[anchor=west,font=\scriptsize]{$x$};
        \draw[-{Latex[length=2mm]},ranblue,thick] (0,0) -- (0,3.2) node[anchor=south,font=\scriptsize]{$y$};
        \fill[coreorange] (2,2.2) circle (3pt);
        \draw[dashed,softgray] (2,2.2) -- (2,0) node[anchor=north,font=\scriptsize]{2};
        \draw[dashed,softgray] (2,2.2) -- (0,2.2) node[anchor=east,font=\scriptsize]{2.2};
        \node[font=\scriptsize,text=softgreen,align=center,anchor=north] at (1.6,-0.65)
          {perpendicular:\\two numbers, no overlap};
      \end{tikzpicture}
    \column{0.48\textwidth}
      \centering
      \begin{tikzpicture}[scale=0.72,transform shape]
        \draw[-{Latex[length=2mm]},ranpurple,thick] (0,0) -- (3.2,0) node[anchor=west,font=\scriptsize]{$a$};
        \draw[-{Latex[length=2mm]},ranpurple,thick] (0,0) -- (2.77,1.6) node[anchor=west,font=\scriptsize]{$b$};
        \fill[coreorange] (2,2.2) circle (3pt);
        \draw[dashed,softgray] (2,2.2) -- (2,0);
        \draw[dashed,softgray] (2,2.2) -- (2.45,1.42);
        \node[font=\scriptsize,text=ranpurple,align=center,anchor=north] at (1.6,-0.65)
          {tilted:\\$a$ and $b$ repeat each other};
      \end{tikzpicture}
  \end{columns}

  \vspace{0.6em}
  {\small In 2-D you need two axes, in 3-D three --- and they must be
   \textbf{perpendicular}. Tilt them and each axis partly repeats what the other
   already said: the two numbers overlap, and reading one leaks into the other.\par}

  \vspace{0.3em}
  \takeaway{A signal is a point in a space with many dimensions.\\To describe it cleanly we need many axes, all perpendicular to each other.}
\end{frame}

\begin{frame}{Step 3: the FFT is a set of perpendicular axes}
  \vspace{0.1em}
  \begin{columns}[T,onlytextwidth]
    \column{0.42\textwidth}
      {\small\textbf{Two dimensions.} The 2-point FFT:\par}
      \vspace{-0.6em}
      {\small\[ \begin{pmatrix} 1 & \phantom{-}1 \\ 1 & -1 \end{pmatrix}
          \begin{pmatrix} a \\ b \end{pmatrix} =
          \begin{pmatrix} a+b \\ a-b \end{pmatrix} \]}
      \vspace{0.15em}
      {\footnotesize Row one: the \textbf{common part}. Row two: the \textbf{difference}.
       The rows are perpendicular --- $1\cdot 1 + 1\cdot(-1) = 0$ --- so the two
       answers do not overlap.\par}
    \column{0.54\textwidth}
      {\small\textbf{Four dimensions.} Same idea; $j=\sqrt{-1}$ lets an axis rotate:\par}
      \vspace{-0.6em}
      {\small\[ \begin{pmatrix} 1 & \phantom{-}1 & \phantom{-}1 & \phantom{-}1 \\
                                1 & -j & -1 & \phantom{-}j \\
                                1 & -1 & \phantom{-}1 & -1 \\
                                1 & \phantom{-}j & -1 & -j \end{pmatrix} \]}
      \vspace{0.15em}
      {\footnotesize Any two rows: multiply entry by entry (flip the sign of $j$ in
       one of them), add --- you get 0. \textbf{Every row is perpendicular to every
       other.} Each row is one frequency.\par}
  \end{columns}

  \vspace{0.35em}
  \takeaway{FFT asks a signal: how much of you lies along each frequency?\\IFFT is the reverse: give it the amounts, it builds the signal.}
\end{frame}

\begin{frame}{Step 4: an OFDM symbol is just an IFFT}
  \vspace{0.2em}
  \centering
  \begin{tikzpicture}[scale=0.84,transform shape,
      pbox/.style={draw=ranblue,fill=blue!5,rounded corners=3pt,minimum width=2.05cm,
        minimum height=0.95cm,align=center,font=\scriptsize}]
    \node[pbox] (bits) at (0,0) {bits\\[-1pt]{\tiny 1011 0010 \ldots}};
    \node[pbox] (qam) at (2.9,0) {16-QAM\\[-1pt]{\tiny 64 points}};
    \node[pbox,draw=coreorange,fill=orange!8] (ifft) at (5.8,0) {IFFT\\[-1pt]{\tiny 64 perpendicular axes}};
    \node[pbox,minimum width=2.6cm] (sym) at (8.9,0) {one OFDM symbol\\[-1pt]{\tiny 64 subcarriers, a burst in time}};
    \foreach \a/\b in {bits/qam,qam/ifft,ifft/sym} \draw[flow] (\a) -- (\b);
  \end{tikzpicture}

  \vspace{0.6em}
  \begin{columns}[T,onlytextwidth]
    \column{0.50\textwidth}
      \small
      \begin{itemize}\itemsep4pt
        \item Put one QAM point on each of 64 perpendicular axes.
        \item The IFFT adds them into one burst of time: a \textbf{symbol}.
        \item Each axis is a \textbf{subcarrier}, a narrow carrier of its own.
        \item Perpendicular $\Rightarrow$ no interference, even packed tight.
      \end{itemize}
    \column{0.46\textwidth}
      \small
      \textbf{How tight?} Spacing $= 1/T$, one over the symbol length. The standard
      picks the number: 5G uses 15, 30, 60 or 120\,kHz; Wi-Fi uses 312.5\,kHz.

      \vspace{0.4em}
      Send symbol after symbol and you get a \textbf{grid}: subcarriers down the
      side, symbols across. \textbf{One cell = one subcarrier for one symbol.}
  \end{columns}

  \vspace{0.2em}
  \takeaway{That grid is the OFDM you already use every day.\\Next slide: the same grid, read as a radar.}
\end{frame}

\begin{frame}{OFDM already carries a radar inside it}
  \vspace{0.3em}
  \begin{columns}[T,onlytextwidth]
    \column{0.44\textwidth}
      \centering
      \begin{tikzpicture}[scale=0.80,transform shape]
        \foreach \x in {0,...,5}
          \foreach \y in {0,...,3}
            \draw[draw=softgray!50,fill=blue!6] (\x*0.62+0.12,\y*0.52+0.10)
              rectangle ++(0.58,0.48);
        \draw[-{Latex[length=2mm]},softgray] (0,0) -- (4.0,0)
          node[anchor=north,font=\scriptsize] {time (symbols)};
        \draw[-{Latex[length=2mm]},softgray] (0,0) -- (0,2.4);
        \node[font=\scriptsize,rotate=90,anchor=south] at (-0.16,1.2) {frequency};
      \end{tikzpicture}
    \column{0.52\textwidth}
      \small
      \begin{itemize}\itemsep5pt
        \item Mature: Wi-Fi, 4G, 5G, 6G.
        \item The grid is already knobs.
        \item A cell serves bits or echoes.
        \item No new modulator needed.
      \end{itemize}
  \end{columns}

  \vspace{0.4em}
  \qa{If OFDM is already good enough, what is left to design?}
     {Which cells go to bits and which go to echoes. Every cell handed to sensing
      is a cell taken from data.}

  \glossline{OFDM = Orthogonal Frequency Division Multiplexing}
\end{frame}
